\documentclass[12pt]{amsart}
\usepackage{amsfonts, amsmath, latexsym, epic, eepic}
\usepackage{url}
\title{Perfect complex functionality}
\author{}

\begin{document}
\newcommand{\RR}{\ensuremath{\mathbb{R}}}
\newcommand{\NN}{\ensuremath{\mathbb{N}}}
\newcommand{\QQ}{\ensuremath{\mathbb{Q}}}
\newcommand{\CC}{\ensuremath{\mathbb{C}}}
\newcommand{\ZZ}{\ensuremath{\mathbb{Z}}}
\newcommand{\TT}{\ensuremath{\mathbb{T}}}
\newtheorem{prop}{Proposition}
\newtheorem{theor}{Theorem}
\newtheorem{cor}{Corollary}
\newtheorem{lem}{Lemma}
\newtheorem{conj}{Conjecture}
\newtheorem{claim}{Claim}
\newtheorem{remark}{Remark}
\newtheorem{definition}{Definition}

\maketitle
\thispagestyle{empty}

We are able to work with perfect complex effectively. For that, we provide a number of functionality that can be used from GAP for users of the program.



\section{Installation}

All the code is accessible from the repository {\bf polyhedral\_common} which is available at
\begin{verbatim}
https://github.com/MathieuDutSik
\end{verbatim}

There are two GAP files that need to be read
\begin{verbatim}
gap> Read("common.g");
gap> Read("access_points.g");
\end{verbatim}
They contain the GAP code that encapsulates the access to the C++ functions.

Here, the C++ functions in question are provided by the program {\bf PERF\_SerialPerfectComputation.cpp}. There are several ways to compile it which are documented in {\bf README.md}

The compiled program has to be made available in the {\bf PATH} or the environment variable {\bf POLYHEDRAL\_COMMON\_SOURCE\_CODE}.



\section{$T$-space chosen}

The $T$-space on which the computation is going to be done has to be chosen:
\begin{verbatim}
gap> desc:=GenerateTspaceDescription_classic(4, false);
\end{verbatim}
Which will select the classic Voronoi space for $\ZZ^4$.

For selecting the $T$-space of imaginary quadratic, you can do
\begin{verbatim}
gap> desc:=GenerateTspaceDescription_imag_quad(3, -3, false);
\end{verbatim}
This will select the $T$-space for imaginary quadratic over the Eisenstein integers.

The second argument specifies whether we want to consider only the well rounded cells, which is usually not the case.



\section{Usage}

The record {\bf desc} created in the previous operation contains the information as well as the place where the cache is going to be stored. The complex is created at the first call and then reused in subsequent calls. This means that the code is slow at first and then much faster later on. If you want to specify the precise place where the cache is stored, please adjust the entry {\bf FileCache} in the record {\bf desc}.


\subsection{Getting generators of the group}

One of the basic output of Voronoi method is that it gives you a generating set of the group.

\begin{verbatim}
gap> desc:=GenerateTspaceDescription_imag_quad(3, -3, false);
gap> ListGen:=PERFCOMP_group_generators(desc);
\end{verbatim}
Note that here the argument {\bf false} is needed because we have to compute the.

\subsection{Information on specific vector family}

One wants to find properties of a cell. There are some commands for that
\begin{verbatim}
gap> desc:=GenerateTspaceDescription_imag_quad(3, -3, false);
gap> EXT:=[[0,0,1, 0,0,0]];
gap> record_dim:=PERFCOMP_dimension(desc, EXT);
gap> is_face:=PERFCOMP_is_face(desc, EXT);
gap> face_index:=PERFCOMP_face_search(desc, EXT);
gap> list_lower_cell:=PERFCOMP_lower_boundary_cell(desc, EXT);
gap> list_upper_cell:=PERFCOMP_upper_boundary_cell(desc, EXT);
\end{verbatim}
The {\bf record\_dim} is the record containing the information on the dimension. The {\bf index} in that record is what matters when calling for a chain or something. index=0 corresponds to the perfect matrices.

{\bf is\_face} is a boolean indicating if EXT is a face or not in the complex.
{\bf face\_index} is a record containing the index, the number of orbit that describes it and the matrix transformation that maps the known representative to that one.

{\bf list\_lower\_cell} is the list of cells contained in EXT. Note that if we select well rounded or not, it will change.
{\bf list\_upper\_cell} is the list of cells containing EXT. This requires EXT to be well rounded since otherwise, we would have an infinity of cells.


\subsection{Computing stabilizers and equivalence}

The computation of the perfect complex with the list of triples allow computing the stabilizer as well as testing equivalence of cells.


\begin{verbatim}
gap> desc:=GenerateTspaceDescription_imag_quad(3, -3, false);
gap> EXT:=[[0,0,1, 0,0,0]];
gap> ListGen:=PERFCOMP_stabilizer(desc, EXT);
\end{verbatim}

for equivalence it works in the same way:

\begin{verbatim}
gap> desc:=GenerateTspaceDescription_imag_quad(3, -3, false);
gap> EXT1:=[[1,0,0, 0,0,0]];
gap> EXT2:=[[0,0,1, 0,0,0]];
gap> result:=PERFCOMP_test_equivalence(desc, EXT1, EXT2);
\end{verbatim}

Here it is important to have {\bf only\_well\_rounded = false} since the cells EXT and EXT2 are indeed not well rounded.




\subsection{Chain boundaries, contracting homotopy and simplification}

A chain is expressed as a vector of entries for example

\begin{verbatim}
gap> chain1:=[rec(iOrb:=1, value:=1, M:=M1),
              rec(iOrb:=2, value:=-1, M:=M2)];
\end{verbatim}
That is we select some group element M1, M2 in the group being considered. iOrb is the index of the orbit. value is the coefficient in front.
The length of the chain can be arbitrary large.



\begin{verbatim}
gap> chain2:=PERFCOMP_chain_simplification(desc, 2, chain1);
gap> chain3:=PERFCOMP_chain_boundary(desc, 2, chain1);
gap> chain4:=PERFCOMP_chain_contracting_homotopy(desc, 2, chain1);
\end{verbatim}
The meaning of those operations should be clear here. The simplification leaves the chain at the same index.
The contracting homotopy returns a chain of one index lower. The boundary increases the index.

Here is a few ideas on how to use those functions:
\begin{enumerate}
\item If you want to test if a chain is zero, then you call the simplification. The chain will be zero if and only if the length is 0.
\item You want to compute the boundary. First make sure that the dimension is higher than 1. Then use the above operation.
\item You want to compute the contracting homotopy. You need to first compute the boundary which has to be zero, and then you can compute the contracting homotopy.
\end{enumerate}





\end{document}
